
\section{\extrahard\ 旁注：Koszul \trComplexN{}与\trHilbert{}\trSyzygy{}定理}

为了展示我们迄今发展的理论已经相当有用，我们就用它来证明\trHilbert{}\trSyzygy{}定理~16.1.3。这反过来也给了我们一个绝好的借口，来引入交换代数中的一件重要工具——Koszul \trComplexN。我们不会用到本节的内容，所以唯有在你兴致来临时才去读它。

\textbf{24.4.1.}\ \emph{Koszul \trComplexN。}\quad 设 $E$ 是秩为 $r$ 的自由 $A$-\trModule。给定一个元素 $s\in E^{\vee}$，即一个 $A$-线性映射 $s\colon E\to A$，与 $s$ 相伴的\textbf{Koszul \trComplexN} $\mathbf{K}_\bullet(s)$ 是链\trComplexN
\[
0 \longrightarrow \Lambda^{r}E \overset{\mathrm{d}_r}{\longrightarrow} \Lambda^{r-1}E \overset{\mathrm{d}_{r-1}}{\longrightarrow} \cdots \overset{\mathrm{d}_2}{\longrightarrow} \Lambda^{1}E \overset{\mathrm{d}_1}{\longrightarrow} A \longrightarrow 0.
\]
微分定义为
\[
\mathrm{d}_k(e_1\wedge\cdots\wedge e_k)=\sum_{i=1}^{k}(-1)^{i+1}\,s(e_i)\,e_1\wedge\cdots\wedge\hat{e}_i\wedge\cdots\wedge e_k,
\]
其中 $e_1,\dots,e_r\in E$（而 $\hat{e}_i$ 一如通常，表示「略去项 $e_i$」）。不难看出这是一个\trComplexN（这也可以由\mbox{练习~24.4.B} 推出）；要看出它是\trWellDefined{}的，则要稍难一些。

为方便起见，在 $E$ 中选取一组基，并设 $s\colon A^{\oplus r}\to A$ 把 $E=A^{\oplus r}$ 的第 $i$ 个基元素送到 $f_i$。在这种情形，我们把该 Koszul \trComplexN{}记作 $\mathbf{K}_\bullet(f_1,\dots,f_r)$。\footnote{原书此处作「and let suppose」，其中「let」字当为赘文（当作「suppose」），径改。——译者注；又，原书下句作「we denote Koszul complex by」，脱定冠词「the」，补足。——译者注}

\textbf{24.4.A.（练习）}\quad 若 $A=k[x_1,\dots,x_r]$，亲手验证：Koszul \trComplexN $\mathbf{K}_\bullet(x_1,\dots,x_r)$ \trExact，唯独在最后一项例外，那里\trCohomology{}为 $k=A/(x_1,\dots,x_r)$。（一种可行的方法：盯住每个\trMonomial{}发生的变化，一如 \S18.3。）

这样，Koszul \trComplexN{}给出了 $k[x_1,\dots,x_r]/(x_1,\dots,x_r)$ 的一个自由分解。（\mbox{练习~24.4.C} 将把此事推广。）

\textbf{24.4.B.（练习：Koszul \trComplexN{}的有用解释）}\quad 证明：Koszul \trComplexN $\mathbf{K}_\bullet(f_1,\dots,f_r)$ 就是\trComplexN
\[
\bigotimes_{i=1}^{r}\Bigl(0 \longrightarrow A \overset{\times f_i}{\longrightarrow} A \longrightarrow 0\Bigr).
\]
先搞清楚这是什么意思！

\textbf{24.4.C.（练习）}\quad 证明：若 $f_1,\dots,f_r$ 是一个\trRegularSequence，则 $\mathbf{K}_\bullet(f_1,\dots,f_r)$ 处处\trExact，唯独在最后一项例外，那里\trCohomology{}为 $A/(f_1,\dots,f_r)$。提示：对 $r$ 作归纳，并用\mbox{练习~24.4.B}。

于是，若 $f_1,\dots,f_r$ 是一个\trRegularSequence，Koszul \trComplexN{}便给出 $A/(f_1,\dots,f_r)$ 的一个长度为 $r$ 的自由分解，称为 $A/(f_1,\dots,f_r)$ 的\textbf{Koszul 分解}（Koszul resolution）。

\textbf{24.4.2.}\ \emph{\trhRemark。}\quad （a）Koszul \trComplexN{}的 $n=2$ 情形，曾在\mbox{定理~9.5.6} 的证明中隐秘登场——它正是那个小小的\trDoubleComplex{} \eqref{eq:double-complex-regular-sequence} 的\trTotalComplex。（b）类似地，也可以对一个 $A$-\trModule{} $M$ 作 Koszul \trComplexN：$\mathbf{K}_\bullet(s,M):=\mathbf{K}_\bullet(s)\otimes_A M$。（c）Koszul \trComplexN{}的\trCohomology{}称为\textbf{Koszul 同调}（Koszul homology）。它所度量的，正是一个序列作为\trRegularSequence{}的失败程度。

\textbf{24.4.3. \trRegular\trLocalRing{}上的\trFinitelyGenerated{}\trModule{}都有有限自由分解。}

为接下来的四道练习，设 $(A,\mathfrak{m})$ 是一个（\trNoether{}的）\trRegular\trLocalRing，维数为 $n$，并设 $f_1,\dots,f_n$ 是 $\mathfrak{m}$ 的一组\trGenerators。

\textbf{24.4.D.（练习）}\quad 证明：对\emph{任何} $A$-\trModule{} $M$，只要 $i>n$ 就有 $\mathrm{Tor}_i^A(M,A/\mathfrak{m})=0$。提示：取 $A/\mathfrak{m}$ 的一个合适的自由分解。

\textbf{24.4.E.（练习）}\quad 若对 $i>b\ge 1$ 都有 $\mathrm{Tor}_i^A(M,A/\mathfrak{m})=0$，而 $\alpha\colon A^{\oplus a}\to M$ 是一个\trSurjectiveMap，证明：对 $i>b-1$ 有 $\mathrm{Tor}_i^A(\ker\alpha,A/\mathfrak{m})=0$。

\textbf{24.4.F.（练习）}\quad 证明：每个\trFinitelyGenerated{}的 $A$-\trModule{} $M$ 都有长度至多为 $n$ 的自由分解。提示：由\mbox{定理~24.3.5}，若 $M$ 是满足 $\mathrm{Tor}_i^A(M,A/\mathfrak{m})=0$ 的\trFinitelyGenerated{} $A$-\trModule，则 $M$ 是自由的。

\textbf{24.4.G.（练习）}\quad 举一个例子，说明这个界无法改进。

\textbf{24.4.4. \mbox{\trHilbert{}\trSyzygy{}定理~16.1.3} 的证明。}

我们现在按照 \S24.4.3 所提示的纲要来证明\trHilbert{}\trSyzygy{}定理。设 $\mathbf{S}_\bullet$ 是分次环 $k[x_0,\dots,x_n]$，带通常的分次，其中 $k$ 是一个\trField。回想 \S15.1.4 中定义的分次模 $\mathbf{S}(e)_\bullet$：$\mathbf{S}(e)_d=\mathbf{S}_{e+d}$，于是 $\widetilde{\mathbf{S}(e)_\bullet}=\OO_{\mathbb{P}^n_k}(e)$。

\textbf{24.4.H.（热身练习）}\quad 依照\mbox{练习~24.4.A} 的 Koszul 分解之指引，描述 $\mathbb{P}^n_A$ 上的一个\trExactSequence：
\begin{multline*}
0 \longrightarrow \OO(-n-1) \longrightarrow \cdots \longrightarrow \OO(-j)^{\oplus\binom{n+1}{j}} \longrightarrow \cdots \longrightarrow \\
\OO(-1)^{\oplus(n+1)} \longrightarrow \OO \longrightarrow 0.
\end{multline*}

设 $\mathbf{M}_\bullet$ 是一个\trFinitelyGenerated{}的分次 $\mathbf{S}_\bullet$-\trModule。选取次数为 $e_1,\dots,e_r$ 的\trHomogeneous{}元素 $a_1,\dots,a_r$，使其生成 $\mathbf{M}_\bullet$，于是映射
\[
\times(a_1,\dots,a_r)\colon \bigoplus_{i=1}^{r}\mathbf{S}_\bullet \longrightarrow \mathbf{M}_\bullet
\]
是一个\trSurjectiveMap。我们把它改写成分次模之间的一个保次数映射：
\[
\times(a_1,\dots,a_r)\colon \bigoplus_{i=1}^{r}\mathbf{S}(-e_i)_\bullet \twoheadrightarrow \mathbf{M}_\bullet.
\]

\textbf{24.4.5. \trhProposition（\mbox{定理~24.3.5} 的类似命题）。}\ ——\ \emph{设 $\mathbf{M}_\bullet$ 是一个\trFinitelyGenerated{}的分次 $\mathbf{S}_\bullet$-\trModule，且 $\mathrm{Tor}_1^{\mathbf{S}_\bullet}(\mathbf{M}_\bullet,k)=0$（这里 $k=\mathbf{S}_\bullet/\mathbf{S}_+$）。则对诸 $d_i$ 的某种选取，$\mathbf{M}_\bullet\cong\bigoplus_{i=1}^{n}\mathbf{S}(-d_i)_\bullet$ 作为分次模。\footnote{原书此处直和上限作「$n$」，而其证明（24.4.6）取「$r$」（如取 $\mathbf{M}_\bullet=\mathbf{S}_\bullet$ 即见按 $n$ 不成立），疑当为「$r$」之误，兹照录原书。——译者注}}

换言之，$\mathbf{M}_\bullet$ 是「分次意义下自由」的。为证明此命题，我们先从\trNakayama{}引理的一个分次变体入手。

\textbf{24.4.I.（练习：分次\trNakayama{}引理）}\quad 设 $\mathbf{C}_\bullet$ 是一个\trFinitelyGenerated{}的 $\mathbf{S}_\bullet$-\trModule，满足 $\mathbf{C}_\bullet\otimes_{\mathbf{S}_\bullet}k=0$。证明：$\mathbf{C}_\bullet=0$。提示：若 $\mathbf{C}_\bullet\neq 0$，看其中次数最小的非零元素。（你不会需要动用「通常的」\trNakayama{}引理。）

\textbf{24.4.6.}\ \emph{\mbox{\trhProposition~24.4.5} 的证明。}\quad 设 $\mathrm{Tor}_1^{\mathbf{S}_\bullet}(\mathbf{M}_\bullet,k)=0$。取 $\mathbf{M}_\bullet$ 的\trHomogeneous{}元素 $y_0,\dots,y_r$\footnote{原书此处作「Let $y_0,\dots,y_r$ be homogeneous elements」，其起点与下文的求和指标 $\bigoplus_{i=1}^{r}$ 不协，疑当为「$y_1,\dots,y_r$」，兹照录原书。——译者注}，使它们的像构成 $\mathbf{M}_\bullet/(x_0,\dots,x_n)=\mathbf{M}_\bullet\otimes_{\mathbf{S}_\bullet}k$ 的一组基。我们有映射 $\alpha\colon\bigoplus_{i=1}^{r}\mathbf{S}(-\deg y_i)_\bullet\to\mathbf{M}_\bullet$，它与 $k$ 作张量后便成为一个\trIsomorphism。令 $\mathbf{Q}_\bullet=\coker\alpha$，对
\[
\bigoplus_{i=1}^{r}\mathbf{S}(-\deg y_i)_\bullet \overset{\alpha}{\longrightarrow} \mathbf{M}_\bullet \longrightarrow \mathbf{Q}_\bullet \longrightarrow 0
\]
施加 $\otimes_{\mathbf{S}_\bullet}k$，便见 $\mathbf{Q}_\bullet\otimes_{\mathbf{S}_\bullet}k=0$，从而由\mbox{练习~24.4.I} 得 $\mathbf{Q}_\bullet=0$，故 $\alpha$ 是\trSurjectiveMap。再令 $\mathbf{K}_\bullet=\ker\alpha$，对
\[
0 \longrightarrow \mathbf{K}_\bullet \longrightarrow \bigoplus_{i=1}^{r}\mathbf{S}(-\deg y_i)_\bullet \overset{\alpha}{\longrightarrow} \mathbf{M}_\bullet \longrightarrow 0
\]
施加 $\otimes_{\mathbf{S}_\bullet}k$，并用假设 $\mathrm{Tor}_1^{\mathbf{S}_\bullet}(\mathbf{M}_\bullet,k)=0$，便见 $\mathbf{K}_\bullet\otimes_{\mathbf{S}_\bullet}k=0$，从而 $\mathbf{K}_\bullet=0$，故 $\alpha$ 给出我们所要的\trIsomorphism。\hfill$\square$

现在，我们已经完全可以照搬 \S24.4.3 的模板了。

\textbf{24.4.J.（练习）}\quad 证明：每个\trFinitelyGenerated{}的分次 $\mathbf{S}_\bullet$-\trModule{} $\mathbf{M}_\bullet$\footnote{原书此处作「every finite-generated graded $\mathbf{S}_\bullet$-module」，「finite-generated」当为「finitely generated」之误，径改。——译者注}都有长度至多为 $n+1$ 的「分次自由分解」：
\[
0 \longrightarrow \mathbf{F}_\bullet^{n+1} \longrightarrow \cdots \longrightarrow \mathbf{F}_\bullet^{1} \longrightarrow \mathbf{M}_\bullet \longrightarrow 0,
\]
其中 $\mathbf{F}_\bullet^q=\bigoplus_{i=1}^{j_q}\mathbf{S}(-a_{i,q})_\bullet$。

\textbf{24.4.K.（练习）}\quad 证明\mbox{\trHilbert{}\trSyzygy{}定理~16.1.3}。
